\noindent\textbf{例 5.10}\quad (加班时数的 $T^{2}$ 控制图)

利用例 5.8 中警署的数据,可以构造两变量 $X_1=$法定加班出勤时数及 $X_2=$特殊事件处理加班时数的 $T^{2}$ 控制图. 变量多于两个的 $T^{2}$ 控制图在练习 5.26 中考虑. 与例 5.9 的椭圆控制图一样,我们取 $\alpha=0.01$.

\hyperref[fig:5.8]{图 5.8} 的 $T^{2}$ 控制图表明,第 11 期的那对数据失控. 与例 5.9 一样,进一步的研究表明,这是由于这段时间内,用于特殊事件处理的加班时数很大所致.

\begin{mvfigure}
\begin{tikzpicture}[x=0.6cm,y=0.32cm]
  \draw[->] (0,0) -- (17,0) node[right,font=\small]{时期};
  \draw[->] (0,0) -- (0,13) node[above,font=\small]{$T^{2}$};
  \foreach \y in {0,2,4,6,8,10,12} \draw (0,\y) -- (-0.18,\y) node[left,font=\scriptsize]{\y};
  \draw[dashed] (0,9.21) -- (16.5,9.21);
  \draw (1,0.5)--(2,2.3)--(3,2.6)--(4,1.0)--(5,1.0)--(6,2.5)--(7,0.8)--(8,2.5)--(9,0.5)--(10,0.6)--(11,10.4)--(12,7.6)--(13,0.7)--(14,0.8)--(15,0.4)--(16,0.5);
  \foreach \x/\y in {1/0.5,2/2.3,3/2.6,4/1.0,5/1.0,6/2.5,7/0.8,8/2.5,9/0.5,10/0.6,11/10.4,12/7.6,13/0.7,14/0.8,15/0.4,16/0.5}
    \fill (\x,\y) circle (1.5pt);
\end{tikzpicture}
\phantomsection\label{fig:5.8}\mvcap{图 5.8\quad 法定加班出勤时数及特殊事件处理加班时数的 $T^{2}$ 质量控制图($\alpha=0.01$)}
\end{mvfigure}

当多元 $T^{2}$ 控制图发出第 $j$ 个数据落在控制域外面的信号时,就应确定哪些变量要对此负责. 为此,我们经常选择基于庞弗罗尼区间的控制域来解决这一问题. 如果 $x_{jk}$ 不落在区间 $\bigl(\bar{x}_k-t_{n-1}(0.005/p)\sqrt{s_{kk}},\ \bar{x}_k+t_{n-1}(0.005/p)\sqrt{s_{kk}}\bigr)$ 之中,则第 $k$ 个变量失控,此处 $p$ 为所测量变量的总数.

\noindent\textbf{例 5.11}\quad (焊接机器的控制——仅有 $T^{2}$ 控制图是不够的)

汽车的传动轴的装配需要进行管状物嵌套的圆形焊接. 自动焊接机器的输入必须控制在一定的操作范围里,这样才能保证焊接质量. 为了控制这个过程,监控工程师观测了四个关键的变量:
\[
X_1=\text{电压(伏特)},\quad X_2=\text{电流(安培)},
\]
\[
X_3=\text{进给速度(英寸/分)},\quad X_4=\text{(输入)气流(立方英尺/分)}
\]
\hyperref[tab:5.9]{表 5.9} 中给出了在 5 秒时间内的这些变量的数据.

\begin{mvtable}
\phantomsection\label{tab:5.9}\mvtabcap{表 5.9\quad 焊接数据}
\begin{tabular}{cccccc}
\toprule
序号 & 电压($X_1$) & 电流($X_2$) & 进给速度($X_3$) & 气流($X_4$)\\
\midrule
1 & 23.0 & 276 & 289.6 & 51.0\\
2 & 22.0 & 281 & 289.0 & 51.7\\
3 & 22.8 & 270 & 288.2 & 51.3\\
4 & 22.1 & 278 & 288.0 & 52.3\\
5 & 22.5 & 275 & 288.0 & 53.0\\
6 & 22.2 & 273 & 288.0 & 51.0\\
7 & 22.0 & 275 & 290.0 & 53.0\\
8 & 22.1 & 268 & 289.0 & 54.0\\
\bottomrule
\end{tabular}
\end{mvtable}
